%! AP Statistics — Unit 2, Lesson 2.4 — Warm-Up
\documentclass[10pt]{article}
\usepackage{apstats-article}
\usepackage{apstats-boxes}

%=============================================================================
\begin{document}

\pageheader{Unit 2, Lesson 2.4}{Warm-Up}
\namedateperiod

\vspace{0.3in}

% ── SECTION 1: Spiral Review ──────────────────────────────────────────────────
{\large\bfseries\color{navy} Part 1 \quad Explanatory and Response Variables (Lesson 2.1 Review)}

\vspace{0.12in}

\noindent For each scenario below, identify the \textbf{explanatory variable} and the
\textbf{response variable}. Then state which axis each belongs on in a scatterplot.

\vspace{0.1in}

\begin{enumerate}[leftmargin=1.6em, itemsep=0.22in]

  \item A researcher wonders whether the number of hours a student studies predicts their
        score on the next exam.

        \vspace{8pt}
        Explanatory variable: \blank{5.5cm} \quad Response variable: \blank{5.5cm}

  \item A nutritionist investigates whether a person's daily caloric intake is associated
        with their body mass index (BMI).

        \vspace{8pt}
        Explanatory variable: \blank{5.5cm} \quad Response variable: \blank{5.5cm}

  \item A city planner records daily high temperature and daily ice cream sales at a local
        shop to look for a pattern.

        \vspace{8pt}
        Explanatory variable: \blank{5.5cm} \quad Response variable: \blank{5.5cm}

\end{enumerate}

\vspace{0.28in}

% ── SECTION 2: Looking Ahead ──────────────────────────────────────────────────
{\large\bfseries\color{navy} Part 2 \quad Setting Up for Scatterplots}

\vspace{0.12in}

\begin{tcolorbox}[colback=greenbg, colframe=greenacc, boxrule=0.8pt, arc=2mm,
    left=4mm, right=4mm, top=3mm, bottom=3mm]
    \small
    \textbf{Think about it:} A scatterplot displays two quantitative variables at once ---
    one on each axis. Before you can draw one, you need to know \emph{which variable belongs
    where} and \emph{how to label the axes} completely (variable name \textbf{and} units).

    \vspace{6pt}
    Two questions to consider before the lesson:
    \begin{itemize}[leftmargin=1.3em, itemsep=4pt]
        \item \textbf{Axes:} Where does each variable go, and why does it matter?
        \item \textbf{Pattern:} What does a ``positive'' relationship look like visually?
              What does a ``negative'' one look like?
    \end{itemize}
\end{tcolorbox}

\vspace{0.16in}

\noindent Using the temperature vs.\ ice cream sales scenario from \#3 above, sketch what
you think the scatterplot might look like. Label both axes.

\vspace{0.12in}

\begin{center}
\begin{tikzpicture}
  \draw[very thin, gray!40] (0,0) grid (8,5);
  \draw[->, thick, navy] (0,0) -- (8.4,0);
  \draw[->, thick, navy] (0,0) -- (0,5.4);
  \node[below=6pt, navy] at (4,0) {\small \blank{5cm}};
  \node[left=6pt, navy, rotate=90] at (0,2.5) {\small \blank{5cm}};
\end{tikzpicture}
\end{center}

\vspace{0.1in}

\noindent Based on your sketch, does the association appear to be \textbf{positive} or
\textbf{negative}? Explain briefly.\\[8pt]
\writeline\\[2pt]\writeline

\end{document}
